problem:  
Let \(M^n \subset \mathbb{S}^{n+1}(1)\) be a compact hypersurface with **parallel mean curvature vector** (hence constant mean curvature \(H > 0\)) and **constant scalar curvature** \(R\). Classical examples such as Clifford hypersurfaces \(S^p(a)\times S^q(b)\subset S^{n+1}(1)\) (with \(a^2+b^2=1,\; p+q=n\)) satisfy these properties and are known to be stable as PMC hypersurfaces.  

**Question:**  
Assume further that \(M\) has **positive sectional curvature** and that its **Ricci tensor is δ‑pinched** in the sense that  
\[
\frac{\mathrm{Ric}(X,X)}{|X|^2} \geq \delta \, \mathrm{Scal}(M)/n \quad \text{for all } X \in TM
\]
for some constant \( \delta > 0\).  
Determine whether there exists a universal constant \( \delta_0 \in (0,1) \) such that, if \( \delta > \delta_0 \), the only such hypersurfaces are *totally umbilical spheres* or *Clifford hypersurfaces*.  
In other words, is there a **curvature pinching rigidity** for PMC, constant‑scalar‑curvature hypersurfaces in \(\mathbb{S}^{n+1}(1)\)?  

Why is it a "good" problem:  
This problem refines known classification results rather than duplicating them. Classical Clifford hypersurfaces exhaust known families of compact PMC constant‑scalar‑curvature examples, but no precise *pinching rigidity* threshold is understood: how positively curved and Ricci‑pinched such a hypersurface can be before it must collapse to a totally umbilical sphere or a standard product. Establishing or disproving such a universal δ‑threshold would unify several classical rigidity phenomena—Simons‑type inequalities, Cheng–Yau curvature pinching, and the geometry of CMC or PMC spheres—into a single quantitative framework.  

The problem remains wide open: even identifying an optimal δ for which rigidity holds is beyond current techniques. It would bridge global Riemannian pinching theory with the extrinsic geometry of submanifolds, and any progress here would illuminate how intrinsic curvature constraints control extrinsic shape—a central, profound theme in modern differential geometry.